\section{¿Por qué se necesita Retrieval Augmentation?}

El curso primero sitúa la retrieval dentro de la historia del language modeling, en lugar de partir del nombre de algún producto de 2023. Así se distinguen tres niveles: el autoregressive objective es antiguo, la escala aporta nuevas capacidades, el instruction tuning corrige la interfaz entre personas y máquinas, y la retrieval se ocupa de la actualización del conocimiento, la procedencia de la evidencia y una memoria controlable.

\subsection{Los modelos de lenguaje no surgieron en los últimos años}

La slide Age of Language Models usa una oración sencilla para ilustrar la next-token factorization: al predecir ``going'', el ``Where are we'' que la precede es el context. La página siguiente muestra un trabajo de neural language modelling de 1991 y recuerda al lector que la combinación de modelos de lenguaje y redes neuronales tiene ya décadas de historia.

\lecturefigure{slide-02-age-of-language-models.jpg}{Un modelo de lenguaje descompone la probabilidad de una secuencia en probabilidades condicionales token por token}{Video de la clase de Stanford 00:02:21}

\begin{importantbox}{Autoregressive factorization}
Para una secuencia de tokens $w_{1:T}$, un modelo de lenguaje autorregresivo se escribe como
\[
p(w_{1:T})=\prod_{t=1}^{T}p(w_t\mid w_{<t}).
\]
$w_{<t}$ es el contexto anterior a la posición $t$; el modelo solo necesita estimar una y otra vez la distribución condicional del siguiente token. Los LLM tienen una cantidad enorme de parámetros y de datos de entrenamiento, pero esta descomposición probabilística en sí no es nueva.
\end{importantbox}

\lecturefigure{slide-03-elman-1991-neural-lm.jpg}{Artículo de neural language modelling de 1991 mostrado en clase}{Video de la clase de Stanford 00:02:57}

\begin{knowledgebox}{Lectura de la figura: la página histórica sirve para calibrar conceptos}
Esta página escaneada no pide memorizar una arquitectura de red antigua, sino que recuerda que el distributed lexicon, el neural sequence processing y el modelado probabilístico del lenguaje existen desde hace mucho. El avance actual proviene de la evolución conjunta de la escala, la optimización, el hardware, los datos y la interface; no se puede reducir toda esa historia a que una empresa «inventó los modelos de lenguaje».
\end{knowledgebox}

\teachervoice{Kiela bromea con que se molestaría si alguien creyera que OpenAI inventó los modelos de lenguaje. Este humor cumple una función didáctica de calibración: la autoregressive idea es antigua, y la novedad de los LLM proviene principalmente de la scale y de la pila de entrenamiento moderna.}

\subsection{Next-token prediction y la «interfaz de usuario rota»}

La slide Next Word Prediction representa el sistema como input--generator--output y señala que uno de los problemas de la plain sequence-to-token es la broken user interface. Antes, los usuarios tenían que construir prompts extraños para que el modelo realizara una tarea; el instruction tuning y el alignment al estilo de ChatGPT permiten expresar la intención directamente.

\lecturefigure{slide-04-next-word-prediction.jpg}{Estructura simple de la plain next-word prediction y su problema de interfaz}{Video de la clase de Stanford 00:04:30}

\begin{knowledgebox}{Lectura de la figura: el objective del modelo difiere del objective del usuario}
El preentrenamiento solo exige predecir el siguiente token del corpus, pero el usuario espera que el modelo siga instrucciones, rechace solicitudes peligrosas, mantenga el formato y complete la tarea. El prompting, el instruction fine-tuning y el RLHF son interface alignment: facilitan invocar el modelo, pero no resuelven automáticamente el conocimiento desactualizado, la falta de evidencia ni los errores factuales.
\end{knowledgebox}

\teachervoice{El ponente resume la contribución clave de ChatGPT como «fix the user interface». A las personas se les da mejor decir directamente lo que quieren que adivinar qué prompt extraño desencadena una conducta; además, los instruction data son muy escasos en el corpus web ordinario.}

\subsection{Una vez corregida la elicitation de salidas, persisten cinco tipos de carencias}

La slide Eliciting Outputs coloca prompt, instruction tuning y alignment antes del generator, y a la derecha enumera hallucination, attribution, staleness, revisions y customization. Estos cinco puntos no son una misma tasa de error: corresponden, respectivamente, a los hechos, las fuentes, el tiempo, la capacidad de actualización y los límites de usuario/organización.

\lecturefigure{slide-05-eliciting-outputs.jpg}{Carencias del modelo que persisten tras mejorar la elicitation de salidas}{Video de la clase de Stanford 00:06:09}

\begin{knowledgebox}{Términos clave: los cinco problemas no deben mezclarse en una sola «exactitud»}
\begin{tabularx}{\textwidth}{L{0.18\textwidth}>{\raggedright\arraybackslash}X >{\raggedright\arraybackslash}X}
\toprule
Problema & Manifestación & Qué puede aportar la retrieval \\
\midrule
Hallucination & La salida contradice fuentes de hechos aceptables & Aporta evidencia explícita, aunque el generador aún puede ignorarla \\
Attribution & No se puede señalar la fuente & Devuelve el ID del documento, el pasaje y la marca de tiempo \\
Staleness & El conocimiento paramétrico está desactualizado & Reemplaza o actualiza el index sin reentrenar todo el modelo \\
Revision & Es difícil corregir con precisión un hecho erróneo & Modifica los documentos o la política de acceso \\
Customization & Cada organización necesita límites de conocimiento distintos & Cambia entre index privados/de dominio y sus permisos \\
\bottomrule
\end{tabularx}
\end{knowledgebox}

\begin{warningbox}{Recuperar evidencia no equivale a grounded generation}
El retriever puede encontrar el documento correcto y aun así el generator puede citar el pasaje equivocado, mezclar memoria paramétrica o ignorar por completo la evidencia. Por ello hay que evaluar por separado, como mínimo, la retrieval quality, la evidence utilization y la final answer correctness.
\end{warningbox}

\teachervoice{El ponente subraya que, sobre todo en escenarios enterprise con strict accuracy requirements, que el modelo «make up stuff» con alta confianza, que no dé attribution y que no pueda revisarse ni personalizarse con rapidez impide su despliegue real.}

\subsection{Contextualization: la memoria externa entra al generador}

La Contextualization architecture añade documents, document encoder, retriever, query encoder y context. La entrada del usuario ya no va solo al generator, sino que también forma una query; el retriever busca evidencia en la colección externa de documentos y entrega los resultados al modelo junto con el prompt.

\lecturefigure{slide-06-contextualization-architecture.jpg}{Arquitectura RAG básica documents--retriever--context--generator}{Video de la clase de Stanford 00:06:48}

\begin{knowledgebox}{Glosario de primer uso: Retriever, Index, Embedding, Context}
\begin{tabularx}{\textwidth}{L{0.17\textwidth}>{\raggedright\arraybackslash}X}
\toprule
Término & Significado \\
\midrule
Document chunk & Unidad recuperable que se corta de un documento; su granularidad determina la exhaustividad y el costo de contexto \\
Embedding & Mapea una query o un chunk a un vector para poder calcular similitudes \\
Index & Estructura de datos que permite búsquedas; puede ser un índice invertido, un índice vectorial o un índice híbrido \\
Retriever & A partir de la query, devuelve desde el index chunks candidatos y sus puntuaciones \\
Context & Texto de evidencia que finalmente se entrega al generator; no equivale al index completo \\
\bottomrule
\end{tabularx}
\end{knowledgebox}

\lecturefigure{slide-07-two-paradigms.jpg}{Dos pares de paradigmas: closed-book/open-book y parametric/semi-parametric}{Video de la clase de Stanford 00:07:48}

\begin{importantbox}{Parametric y non-parametric memory}
La parametric memory es el conocimiento estadístico que se forma en los pesos del modelo mediante el entrenamiento; actualizarla suele requerir continuar el entrenamiento. La non-parametric memory está formada por documentos, bases de datos, índices y cachés que pueden leerse, escribirse o reemplazarse directamente. A RAG se le suele llamar semi-parametric porque el generador tiene parámetros, mientras que la memory externa puede cambiar en tiempo de prueba.
\end{importantbox}

\subsection{Por qué un external index puede mitigar parte de los problemas}

La slide Why does this solve the issues ofrece dos beneficios directos: reemplazar el index permite la customization, evita la staleness y corrige errores; el grounding puede reducir la hallucination y aportar citation y attribution. Obsérvese que la palabra es «less» y no «zero».

\lecturefigure{slide-08-why-retrieval-solves-issues.jpg}{Aportes del index externo a la personalización, la actualización y el grounding}{Video de la clase de Stanford 00:08:57}

\begin{knowledgebox}{Lectura de la figura: la capacidad de actualización proviene de cambiar dónde reside el estado}
Si un hecho existe solo en los pesos, corregirlo requiere entrenar y volver a desplegar; si el hecho está en el index, la ruta de actualización se convierte en document update $\rightarrow$ re-index $\rightarrow$ la siguiente retrieval. El costo es que el sistema debe gestionar las versiones de los documentos, los permisos, el chunking y la consistencia de la recuperación.
\end{knowledgebox}

\begin{warningbox}{Una base de conocimiento externa no es automáticamente una fuente de verdad}
El index puede contener contenido erróneo, contradictorio, desactualizado o malicioso. RAG transforma en parte la pregunta «¿el modelo recuerda mal?» en «¿en qué fuentes de datos confía el sistema y cómo las ordena y cita?»; aumenta la controlabilidad, pero convierte la gobernanza de datos y la source selection en responsabilidades explícitas.
\end{warningbox}

\teachervoice{La clase advierte que el index puede reemplazarse, los hechos pueden revisarse y las fuentes pueden citarse, pero el generator aún puede desatender el context. Lo que la aumentación por recuperación exige de verdad es una interfaz eficaz entre el retriever y el generator.}

\subsection{Un solo diagrama de arquitectura genera muchas preguntas de investigación}

La slide Many Questions plantea preguntas en torno al diagrama básico: cómo aprende, escala y preprocesa el sistema, cómo codificar y recuperar documentos, y cómo diseñar el prompt y verificar la respuesta. Así convierte RAG de un trick aislado en un espacio completo de diseño de sistemas y prepara la clasificación posterior por train/test, representation y fusion.

\lecturefigure{slide-09-many-questions.jpg}{Preguntas de aprendizaje, escalamiento, recuperación y verificación en torno a la arquitectura RAG básica}{Video de la clase de Stanford 00:09:57}

\begin{importantbox}{Los seis ejes de diseño de RAG}
\begin{enumerate}
  \item \textbf{Corpus}: qué documentos son accesibles y quién mantiene sus versiones y permisos;
  \item \textbf{Chunking}: cómo dividir, solapar y conservar la jerarquía;
  \item \textbf{Representation}: sparse, dense, late interaction o hybrid;
  \item \textbf{Retrieval policy}: cuándo recuperar, cuántos elementos tomar y si hacerlo en varios saltos;
  \item \textbf{Fusion}: concatenar en el prompt, cross-attention, decoder fusion o token-level memory;
  \item \textbf{Learning/evaluation}: qué componentes reciben gradientes y cómo distinguir los errores de recuperación de los de generación.
\end{enumerate}
\end{importantbox}

\subsection{Resumen de la sección}

La retrieval augmentation aporta principalmente una memoria externa actualizable, atribuible y personalizable. No cambia el next-token objective ni garantiza automáticamente que la generación sea fiel; por eso hay que diseñar el corpus, el index, el retriever, la fusion, el generator y la evaluation como un solo sistema.
